\textbf{(CAP-COLUNI 2018 -- Q18)} Uma forma de se estimar quantas pessoas há em uma passeata é a de considerar que em cada metro quadrado cabem, em média, seis pessoas.
\textit{(Disponível em: www.somatematica.com.br/curiosidades/c101.php. Acesso em: 07 set. 2017.)}

Durante uma passeata na Avenida Aquarianos, as imagens aéreas mostraram que, em certo momento, toda sua extensão e largura estavam ocupados por pessoas. Veja um esboço da avenida vista de cima.

A avenida possui 0,6 km de extensão, pode ser vista como a união do retângulo $ABCD$ e do trapézio $CDEF$ e o ângulo $\widehat{BCF}$ é igual a $135^\circ$. Usando $\sqrt2=1{,}4$, é CORRETO afirmar que o número aproximado de pessoas presentes nessa passeata foi:

\includegraphics{./figs/2018-Q18.png}

\begin{enumerate}
	\item 58.400
	\item 64.800
	\item 74.800
	\item 88.400
\end{enumerate}
